% !PRINTED 349
\weberpage{349}
%<<CONT
must be even — equal to 6. Two real or two conjugate imaginary roots
$x_1, x_2$ always belong to a triple whose third root is real; this
follows from the equation:

\[
x_3 = \Theta (x_1, x_2) = \Theta (x_2, x_1).
\]

We shall accordingly call a triple which contains two real or two conjugate
imaginary roots a \gesperrt{real triple}. If any imaginary roots occur in a
triple, then the conjugate imaginary roots also form a triple, and two such
triples shall be called \gesperrt{conjugate imaginary triples}.

There are then three possibilities:

1. \gesperrt{All real roots.}

2. \gesperrt{One real and eight imaginary roots.}

In this case the four pairs of conjugate imaginary roots must determine four
triples, all of which contain the same real root as third. If we denote
this real root by $(0,0)$, then

\[
(1,2)\ (2,1),\quad (1,1)\ (2,2),\quad (1,0)\ (2,0),\quad (0,1)\ (0,2)
\]

must be the four pairs of conjugate roots.

3. \gesperrt{Three real and six imaginary roots.}

The three real roots must here form a triple. If we take for the real
roots

\weq{1}{(0,0)\ (1,1)\ (2,2),}

then in the two rows

\weq{2}{\begin{aligned} &(1,2)\ (0,2)\ (1,0) \\ &(2,1)\ (2,0)\ (0,1) \end{aligned}}

conjugate imaginary roots cannot occur in the same row, because otherwise
the third root of the row in question would be real.

It is to be shown that the three real triples determined by the conjugate
pairs must together contain all three real points.

For suppose that, among the three triples

\[
(0,0)\ (1,2)\ (2,1),\quad (0,0)\ (2,0)\ (1,0),\quad (0,0)\ (0,1)\ (0,2)
\]

two real ones occur; then the third must also be real, i.\,e.

\[
(1,2)\ (2,1),\quad (2,0)\ (1,0),\quad (0,1)\ (0,2)
\]
%CONT>>
